\documentclass[10pt,a4paper]{amsart}
\usepackage{amsmath,amssymb,amsthm,graphicx}
\usepackage{verbatim}
\newenvironment{reference}{\comment}{\endcomment}
%\newenvironment{reference}{}{}
\newenvironment{slogan}{\comment}{\endcomment}
\newenvironment{history}{\comment}{\endcomment}

% For commutative diagrams we use Xy-pic
\usepackage[all]{xy}

% We use 2cell for 2-commutative diagrams.
\xyoption{2cell}
% Theorem environments.
%
\theoremstyle{plain}
\newtheorem{theorem}[subsection]{Theorem}
\newtheorem{proposition}[subsection]{Proposition}
\newtheorem{lemma}[subsection]{Lemma}

\theoremstyle{definition}
\newtheorem{definition}[subsection]{Definition}
\newtheorem{example}[subsection]{Example}
\newtheorem{exercise}[subsection]{Exercise}
\newtheorem{situation}[subsection]{Situation}

\theoremstyle{remark}
\newtheorem{remark}[subsection]{Remark}
\newtheorem{remarks}[subsection]{Remarks}

\numberwithin{equation}{subsection}

% Macros
%
\def\lim{\mathop{\mathrm{lim}}\nolimits}
\def\colim{\mathop{\mathrm{colim}}\nolimits}
\def\Spec{\mathop{\mathrm{Spec}}}
\def\Hom{\mathop{\mathrm{Hom}}\nolimits}
\def\Ext{\mathop{\mathrm{Ext}}\nolimits}
\def\SheafHom{\mathop{\mathcal{H}\!\mathit{om}}\nolimits}
\def\SheafExt{\mathop{\mathcal{E}\!\mathit{xt}}\nolimits}
\def\Sch{\mathit{Sch}}
\def\Mor{\mathop{\mathrm{Mor}}\nolimits}
\def\Ob{\mathop{\mathrm{Ob}}\nolimits}
\def\Sh{\mathop{\mathit{Sh}}\nolimits}
\def\NL{\mathop{N\!L}\nolimits}
\def\CH{\mathop{\mathrm{CH}}\nolimits}
\def\proetale{{pro\text{-}\acute{e}tale}}
\def\etale{{\acute{e}tale}}
\def\QCoh{\mathit{QCoh}}
\def\Ker{\mathop{\mathrm{Ker}}}
\def\Im{\mathop{\mathrm{Im}}}
\def\Coker{\mathop{\mathrm{Coker}}}
\def\Coim{\mathop{\mathrm{Coim}}}

% Boxtimes
%
\DeclareMathSymbol{\boxtimes}{\mathbin}{AMSa}{"02}

%
% Macros for moduli stacks/spaces
%
\def\QCohstack{\mathcal{QC}\!\mathit{oh}}
\def\Cohstack{\mathcal{C}\!\mathit{oh}}
\def\Spacesstack{\mathcal{S}\!\mathit{paces}}
\def\Quotfunctor{\mathrm{Quot}}
\def\Hilbfunctor{\mathrm{Hilb}}
\def\Curvesstack{\mathcal{C}\!\mathit{urves}}
\def\Polarizedstack{\mathcal{P}\!\mathit{olarized}}
\def\Complexesstack{\mathcal{C}\!\mathit{omplexes}}
% \Pic is the operator that assigns to X its picard group, usage \Pic(X)
% \Picardstack_{X/B} denotes the Picard stack of X over B
% \Picardfunctor_{X/B} denotes the Picard functor of X over B
\def\Pic{\mathop{\mathrm{Pic}}\nolimits}
\def\Picardstack{\mathcal{P}\!\mathit{ic}}
\def\Picardfunctor{\mathrm{Pic}}
\def\Deformationcategory{\mathcal{D}\!\mathit{ef}}
\usepackage{cleveref}
\crefname{originalsection}{original section}{original sections}
\Crefname{originalsection}{Original section}{Original sections}
\makeatletter
\newlabel{section-limits-colimits-PSh}{{00VB}{source}{}{sourcepointer-1}{}}
\newlabel{section-limits-colimits-PSh@cref}{{[originalsection][0][]00VB}{[1][0][]source}{}{}{}}
\makeatother
\crefname{originallemma}{original lemma}{original lemmas}
\Crefname{originallemma}{Original lemma}{Original lemmas}
\makeatletter
\newlabel{lemma-colimit-sheaves}{{00WI}{source}{}{sourcepointer-2}{}}
\newlabel{lemma-colimit-sheaves@cref}{{[originallemma][0][]00WI}{[1][0][]source}{}{}{}}
\makeatother
\title{Sections of a colimit of sheaves indexed by an ordinal of large cofinality}
\author{The Stacks Project Authors}
\date{stacks-project:sites.tex:local-snapshot-20261012; GNU FDL 1.2 or later}
\begin{document}
\maketitle
\noindent\textit{Editorial note.} All delivered mathematical text is the author's own, copied from the pinned Stacks Project chapter file sites-cohomology.tex; no statement or proof step was rewritten, reordered or removed.\par\medskip
\noindent\textit{Editorial note.} The proof is the author's complete argument; every step the author writes out is inside the counted span and no closing sentence is a placeholder.\par\medskip
\noindent\textit{Editorial note.} Every labelled reference occurring inside the counted spans is carried in assembly.external\_references as a non-load-bearing navigation pointer (same-chapter ones at their definition site, cross-chapter ones at the invocation inside the counted span), so no load-bearing mathematics is replaced by a stub. No bibliography citation occurs inside the counted spans.\par\medskip
\noindent\textit{Editorial note.} The checked chapter list for same-name collisions is the set of stacks chapters staged in this repository: chow, cohomology, divisors, etale-cohomology, flat, more-algebra, more-morphisms, morphisms, spaces-more-morphisms, varieties together with this card's chapter sites-cohomology.\par\medskip
\noindent\textit{Editorial note.} The preamble is the chapter's own retained minimal preamble (sources/stacks-sites-cohomology/retained-preamble.tex), a verbatim excerpt of the native preamble: the Xy-pic setup and the native theorem declarations with the mathematical macros.\par\medskip
\noindent\textit{Allowed standard prerequisites (editorial).}\begin{itemize}
\item Colimits of sheaves as sheafified colimits of presheaves. Same chapter, Lemma lemma-colimit-sheaves. Identifies the colimit of the system of sheaves with the sheafification of the colimit computed in presheaves.
\item The plus construction and sheafification. Same chapter, the description of the sheafification as the result of applying the plus construction twice. Supplies the criterion for a presheaf to be a sheaf used at the end of the proof.
\end{itemize}\medskip
\ifdefined\hypertarget\hypertarget{sourcepointer-1}{}\fi
\noindent\textit{Source reference.} Original section 00VB: sources/stacks-sites/pinned-original-source.tex, source lines 3938--3938. Navigation pointer invoked at proof line 3938; the counted span contains the pointer only and the cited result is not reproduced inside it, so the pointer is registered as a non-load-bearing navigation stub in the established form for Stacks material. No bibliography citation is involved.\par\medskip
\ifdefined\hypertarget\hypertarget{sourcepointer-2}{}\fi
\noindent\textit{Source reference.} Original lemma 00WI: sources/stacks-sites/pinned-original-source.tex, source lines 3941--3941. Navigation pointer invoked at proof line 3941; the counted span contains the pointer only and the cited result is not reproduced inside it, so the pointer is registered as a non-load-bearing navigation stub in the established form for Stacks material. No bibliography citation is involved.\par\medskip
\begin{lemma}
\label{lemma-colimit-over-ordinal-sections}
Let $\mathcal{C}$ be a site. Let $\beta$ be an ordinal.
Let $\beta \to \Sh(\mathcal{C})$, $\alpha \mapsto \mathcal{F}_\alpha$
be a system of sheaves over $\beta$. For $U \in \Ob(\mathcal{C})$
consider the canonical map
$$
\colim_{\alpha < \beta} \mathcal{F}_\alpha(U)
\longrightarrow
\left(\colim_{\alpha < \beta} \mathcal{F}_\alpha \right)(U)
$$
If the cofinality of $\beta$ is large enough, then this map
is bijective for all $U$.
\end{lemma}

\begin{proof}
The left hand side is the value on $U$ of the colimit $\mathcal{F}_{\colim}$
taken in the category of presheaves, see
Section \ref{section-limits-colimits-PSh}.
Recall that $\colim_{\alpha < \beta} \mathcal{F}_\alpha$ is the
sheafification $\mathcal{F}_{\colim}^\#$ of $\mathcal{F}_{\colim}$, see
Lemma \ref{lemma-colimit-sheaves}.
Let $\mathcal{U} = \{U_i \to U\}_{i \in I}$ be an element of the
set $\text{Cov}(\mathcal{C})$ of coverings of $\mathcal{C}$.
If the cofinality of $\beta$ is larger than the
cardinality of $I$, then we claim
$$
H^0(\mathcal{U}, \mathcal{F}_{\colim}) =
\colim H^0(\mathcal{U}, \mathcal{F}_\alpha) =
\colim \mathcal{F}_\alpha(U) = \mathcal{F}_{\colim}(U)
$$
The second and third equality signs are clear. For the first, say
$s = (s_i) \in H^0(\mathcal{U}, \mathcal{F}_{\colim})$.
Then for each $i$ the element $s_i$ comes from an element
$s_{i, \alpha_i} \in \mathcal{F}_{\alpha_i}(U_i)$ for some
$\alpha_i < \beta$. By the assumption on cofinality, we can choose
$\alpha_i = \alpha$ independent of $i$. Then $s_i$ and $s_j$
map to the same element of $\mathcal{F}_{\alpha_{i, j}}(U_i \times_U U_j)$
for some $\alpha_{i, j} < \beta$. Since the cardinality if
$I \times I$ is also less than the cofinality of $\beta$, we see
that we may after increasing $\alpha$ assume $\alpha_{i, j} = \alpha$
for all $i, j$. This proves that the natural map
$\colim H^0(\mathcal{U}, \mathcal{F}_\alpha)
\to H^0(\mathcal{U}, \mathcal{F}_{\colim})$ is surjective.
A very similar argument shows that it is injective.
In particular, we see that $\mathcal{F}_{\colim}$ satisfies
the sheaf condition for $\mathcal{U}$.
Thus if the cofinality of $\beta$ is larger than the
supremum of the cardinalities of the set of index sets $I$
of coverings, then we conclude.
\end{proof}

\end{document}
